\section{Research Questions}
\label{sec:intro:rqs}

%TODO eq:prelim:role, eq:prelim:inversion, and eq:prelim:retention await chapter 2

\subsection{Problem Formalization}
\label{sec:intro:formalization}
Let $G$ be an AIG with the node roles
of~\eqref{eq:prelim:role} and edge inversions of~\eqref{eq:prelim:inversion} as its features,
and let $S$ be a deterministic synthesis algorithm
(\ref{sec:prelim:algorithms:orchestrate}) mapping $G$ to the function-equivalent,
optimized $S(G)$. Primary inputs, primary outputs, and the constant are invariant under $S$,
so the AND nodes $\mathcal{V}_{\wedge}$ are the only ones it can remove. The
\emph{optimizability} of $G$ under $S$ is the fraction of AND nodes the algorithm removes,
\begin{equation}
    y(G) = 1 - \frac{\left| \mathcal{V}_{\wedge}(S(G)) \right|}{\left| \mathcal{V}_{\wedge}(G) \right|}
    \;\in\; [0, 1].
    \label{eq:intro:target}
\end{equation}
Evaluating~\eqref{eq:intro:target} requires executing $S$, a cost that prediction avoids. An
encoder $f_\theta$, a GNN in the sense of~\ref{sec:prelim:gnn:mp} with a graph-level readout,
is fitted to approximate $y$. A reduction
operator $R$ (\ref{sec:prelim:reduction}) acts ahead of the encoder, so the argument $f_\theta$
receives is $G$ or $R(G)$. Accuracy is reported as predictive error (Smooth L1 loss, RMSE,
and $R^2$) and as circuit-ranking ability (Spearman correlation), both defined
in~\ref{sec:method:experiment:metrics:accuracy} and~\ref{sec:method:experiment:metrics:ranking}.

%TODO compression is still the metric name in 1.3, 2.5, and 4.4
\emph{Node and edge retention} is reported separately as the ratios
$\eta_{\mathcal{V}}$ and $\eta_{\mathcal{E}}$ of~\eqref{eq:prelim:retention}. Sparsification
lowers $\eta_{\mathcal{E}}$ with $\eta_{\mathcal{V}} = 1$, so the two are not
interchangeable. Memory and time are compared per graph against the full-graph
baseline (\ref{sec:method:experiment:metrics:stats}).

\subsection{RQ1: Task Feasibility and Baseline}
\label{sec:intro:rq1}
What predictive accuracy (Smooth L1 loss, RMSE, and $R^2$) and circuit-ranking ability
(Spearman correlation) does a GNN reach on optimizability predicted directly from a
circuit's full, unreduced AIG representation, relative to trivial predictors and
published circuit-representation models?

The RQ1 baseline withholds whole designs (design-disjoint), whereas conventional protocols
often withhold individual graphs (random) or whole synthesis scripts (script-disjoint). By
how much does the RQ1 result change under each of those two protocols
(\ref{sec:method:experiment:splitting},~\ref{sec:results:rq1:protocol})?

\subsection{RQ2: Reduction Efficiency}
\label{sec:intro:rq2}
Applied to AIGs, how do partitioning, sparsification, and summarization/coarsening differ in
node and edge retention, peak memory, training and inference throughput, and the offline cost
of the reduction itself?

\subsection{RQ3: Predictive Retention and Trade-offs}
\label{sec:intro:rq3}
To what extent do models trained on reduced AIGs retain predictive accuracy (Smooth L1 loss,
RMSE, and $R^2$) and circuit-ranking ability (Spearman correlation) relative to the full-graph
baseline of RQ1, and how does the accuracy lost trade off against the memory and time saved?

Does the provably lossless coarsening of \citet{bollen2023exact} apply to AIGs, and what
node and edge retention does it achieve with and without topological level in the initial
coloring? How do inexact summarization methods compare against that lossless reference
point in node and edge retention, and in accuracy?

\subsection{RQ4: Value of Domain-Informed Adaptation}
\label{sec:intro:rq4}
Do domain-informed reductions, which exploit topological level, fanout structure, or
gate role, retain more predictive accuracy than their generic counterparts at matched node
and edge retention, $\eta_{\mathcal{V}}$ and $\eta_{\mathcal{E}}$?

\subsection{RQ5: Cross-State Inference}
\label{sec:intro:rq5}
Can a model trained exclusively on reduced AIGs, whether partitioned, sparsified, or
summarized, be queried on full ones at inference time without falling below its own
accuracy on reduced inputs?

A \emph{structural state} is the form in which a graph reaches the encoder, the reduced
$R(G)$ or the unreduced $G$ of~\ref{sec:intro:formalization}. Cross-state inference
evaluates weights trained in one state on inputs in the other.

\subsection{Scope and Delimitations}
\label{sec:intro:rqs:scope}
Optimizability prediction is the evaluation task, not the object of study. The delimitations
below were fixed in advance. Limitations found during the work are reported
in~\ref{sec:discussion:limitations}.

\paragraph{Target synthesis algorithm.}
\label{sec:intro:scope:algorithm}
Training and evaluation target one synthesis algorithm, Orchestrate
(\ref{sec:prelim:algorithms:orchestrate}), since variance across algorithms would confound
the cost attributable to a reduction.

\paragraph{Graph scale.}
\label{sec:intro:scope:scale}
RQ1's full-graph baseline must fit in memory, so the corpus is bounded at an intermediate
rather than an industrial scale
(\ref{sec:method:experiment:scale},~\ref{sec:discussion:limitations:scalability}).

\paragraph{Definition of optimizability.}
\label{sec:intro:scope:target}
Optimizability is relative \emph{node} reduction only. Area, delay, and power, the objectives
synthesis actually optimizes for, are not modeled
(\ref{sec:discussion:limitations:validity}).

\paragraph{Reduction families.}
\label{sec:intro:scope:reductions}
Three reduction families are in scope: partitioning, sparsification, and
summarization/coarsening. This set is the taxonomy of \citet{hashemi2024survey} with
partitioning in place of condensation. Condensation synthesizes a new small graph whose nodes
correspond to no nodes of the original, and RQ5 requires that correspondence
(\ref{sec:relwork:reduction:summarization}).

\paragraph{Model architecture.}
\label{sec:intro:scope:architecture}
One encoder, an edge-aware GCN+ (\ref{sec:method:architecture}), is held fixed so that
outcome differences are attributable to the reduction.

\paragraph{Prediction rather than script generation.}
\label{sec:intro:scope:prediction}
Optimizability is predicted, not synthesis scripts. The prediction is bound to one algorithm:
it answers whether that algorithm is worth running on a given circuit, and by how much.
Selecting a script instead requires ranking \emph{algorithms} for one circuit, which a
single-algorithm model cannot support (\ref{sec:conclusion:futurework:generation}).

